% from ega4-1-fr.tex
\oldpage[IV-1]{236}

(1.6.1), il faut et il suffit que $f'$ et $f''$ la possèdent~; c'est évident
pour la propriété d'être localement de présentation finie, qui est locale sur
$X$~; pour la propriété d'être quasi-compact, cela a été vu en (1.1.6) et pour
la propriété d'être quasi-séparé, en (1.2.9).
\end{proof}

\subsection{Améliorations de résultats antérieurs}
\label{subsection:IV.1.7-fr}

Nous allons donner dans ce numéro une liste de propositions démontrées dans
les chapitres précédents, dont l'énoncé peut être amélioré à l'aide des
nouvelles conditions de finitude introduites ci-dessus.

\begin{env}[1.7.1]
\label{IV.1.7.1-fr}
Dans les énoncés de (I, 6.4.2, 6.4.3 et 6.4.9), on peut, au lieu de supposer
que $X$ et $Y$ sont de type fini sur le corps $K$, supposer seulement qu'ils
sont \emph{localement de type fini} sur $K$~; pour (I, 6.4.2), il suffit
d'observer que tout point d'un préschéma $X$ qui est fermé dans tout ouvert
affine de $X$ est fermé dans $X$, et un schéma affine localement de type fini
sur $K$ est \emph{ipso facto} de type fini (1.5.1). Pour (I, 6.4.9), on utilise
(1.3.7) et (1.3.4, (v)).
\end{env}

\begin{env}[1.7.2]
\label{IV.1.7.2-fr}
Dans (I, 6.5.1, (ii)), on peut au lieu de supposer que $S$ est localement
noethérien et $Y$ de type fini sur $S$, supposer seulement que $Y$ est
\emph{localement de présentation finie} sur $S$, comme le montre aussitôt la
démonstration (en appliquant la définition (1.4.2)). De même dans (I, 6.5.4,
(ii)) et (I, 6.5.5, (ii)) il suffit de supposer que $f$ est un $S$-morphisme,
$Y$ étant localement de présentation finie sur $S$.
\end{env}

\begin{env}[1.7.3]
\label{IV.1.7.3-fr}
Dans (I, 7.1.11), pour la deuxième assertion, au lieu de supposer $S$
localement noethérien et $Y$ de type fini sur $S$, on peut seulement supposer
$Y$ localement de présentation finie sur $S$ (la démonstration dépendant de
(I, 6.5.1, (ii)))~; de même dans (I, 7.1.12 à 7.1.15).
\end{env}

\begin{env}[1.7.4]
\EGAReaderTarget{ega4.c04.statement.1-7-4}{IV.1.7.4-fr}
Dans (I, 9.2.2), on peut remplacer les trois hypothèses par la seule hypothèse
(entraînée par chacune des autres) que $f$ est \emph{quasi-compact et
quasi-séparé}, en vertu de (1.2.6) et (1.2.7). On notera que dans (I, 9.2.1),
l'hypothèse signifie exactement que $f$ est quasi-compact et quasi-séparé.
\end{env}

\begin{env}[1.7.5]
\label{IV.1.7.5-fr}
Dans (I, 9.3.1, 9.3.2 et 9.3.3), on peut remplacer les deux hypothèses sur $X$
par l'hypothèse (moins forte) que $X$ est un préschéma \emph{quasi-compact et
quasi-séparé}, la démonstration utilisant exactement ces deux propriétés (en
vertu de (1.2.7)).
\end{env}

\begin{env}[1.7.6]
\label{IV.1.7.6-fr}
Dans (I, 9.3.4 et 9.3.5), il suffit de supposer $X$ quasi-compact,
$\mathcal F$ et $\mathcal J$ quasi-cohérents et de type fini.
\end{env}

\begin{env}[1.7.7]
\label{IV.1.7.7-fr}
Dans (I, 9.4.7), on peut affaiblir l'hypothèse b) en supposant seulement $X$
\emph{quasi-séparé}, puisqu'il suffit seulement que l'immersion canonique
$U\to X$ soit quasi-compacte (1.2.7). On peut faire la même modification dans
(I, 9.4.8, 9.4.9 et 9.4.10).
\end{env}

\begin{env}[1.7.8]
\label{IV.1.7.8-fr}
Dans (I, 9.5.1), les conditions indiquées entre parenthèses dans l'énoncé
peuvent être remplacées par la condition que $f$ soit quasi-compact et
quasi-séparé.
\end{env}

\begin{env}[1.7.9]
\label{IV.1.7.9-fr}
Dans (I, 9.6.5), on peut remplacer les hypothèses sur $X$ par l'hypothèse
moins forte que $X$ est quasi-compact et quasi-séparé, comme le montre la
démonstration,

